% Page 24 — Alban – XIIIe siècle

\begin{center}
  {\Large\itshape Alban – XIII\textsuperscript{ème} siècle}
\end{center}
\vspace{1.2em}

\noindent
\lettrine{P}{ar ailleurs}, nous savons que le serf n’est pas un esclave. Il a une famille, un foyer, un champ, il n’est pas soumis à un maître, il s’organise comme il veut, et dès lors qu’il a fourni ses redevances, il est quitte envers son Seigneur. Mais, et c’est là une restriction certaine à sa liberté, il est attaché à un domaine. De ce fait, il n’a pas d’obligations militaires quoi qu’il arrive, on ne peut lui arracher sa propriété. Le tenancier libre – lui – peut se voir déposséder de ses biens en cas de dette – et être tenu à combattre s’il y a guerre. Quand au XIV\textsuperscript{ème} siècle les autorités publiques voudront mettre fin au servage, certains intéressés s’y refuseront. C’est ce qu’on appela : « les serfs récalcitrants ».

\bigskip

\noindent
\lettrine{L}{e régime} de la propriété paysanne, consolidé au Moyen Age par le droit d’aînesse, a permis à la famille rurale de se perpétuer sur le même sol pendant des siècles. Le cadastre de 1598 signale déjà des « Cadars » au Noyer, des Villeneuve au Frayssé. Les archives du XVIII\textsuperscript{ème} siècle mentionnent des Sandral et des Trouillet à Margou, des Augé au Verdier. La famille Routoulp qui donna deux curés à Notre Dame d’Alban en 1746 et en 1763 est signalée dans un acte de 1712 comme habitant le masage de « Jouven » où elle se trouve encore. Les exemples abondent, de patients chercheurs pourraient conclure en consultant les registres de baptême ou mariage que les habitants d’Alban sont, mis à part ceux qui sont arrivés récemment dans la région, tous cousins à des degrés divers. Je pense en fournir la preuve dès que possible, en me servant de la liste de mes propres ancêtres originaires par ma grand-mère paternelle du hameau d’Ourtiguet. Comment ne pas s’aimer, quand un même sang coule dans nos veines ?
